\begin{center}
\begin{tikzpicture}[x=0.75cm,y=0.75cm, font=\scriptsize,>=Stealth]
\node[font=\bfseries,anchor=west] at (0,3.3) {火星尘粒：研究测得的尺度};
\draw[->,newsgray] (0.3,.4) -- (8.8,.4) node[right] {粒径};
\foreach \x/\lbl in {1/0.1,3/1,5/10,7/100}{\draw (\x,.3)--(\x,.5);\node[below] at (\x,.3) {\lbl $\mu$m};}
\fill[orange!70] (1,1) circle (.09); \node[anchor=west] at (1.35,1) {超细尘：约$0.1\,\mu$m};
\fill[newsblue!65] (3.8,1.8) circle (.23); \node[anchor=west] at (4.2,1.8) {细尘：低于$2.5\,\mu$m};
\fill[gray!55] (7.0,2.5) circle (.40); \node[anchor=east] at (6.5,2.5) {人类头发约$70\,\mu$m};
\end{tikzpicture}
\par\vspace{1mm}{\scriptsize\color{newsgray}图2｜对数位置仅用于数量级比较，圆面积不按比例。数据：NASA/JPL 2026-10-26新闻稿及VanBommel等PNAS论文；本刊原创示意，非显微照片。}\end{center}